\documentclass[11pt,openany]{book}
\usepackage[utf8]{inputenc}
\usepackage[T1]{fontenc}
\usepackage{lmodern}
\usepackage{amsmath,amssymb,amsthm}
\usepackage{booktabs}
\usepackage{longtable}
\usepackage{listings}
\usepackage{hyperref}
\usepackage[margin=1in]{geometry}
\usepackage{xcolor}
\usepackage{enumitem}
\usepackage{fancyvrb}
\usepackage{microtype}
\usepackage{parskip}

\hypersetup{
  colorlinks=true,
  linkcolor=black,
  urlcolor=blue!60!black,
  citecolor=black,
}

\definecolor{codebg}{RGB}{248,248,248}
\lstset{
  basicstyle=\ttfamily\small,
  backgroundcolor=\color{codebg},
  frame=single,
  breaklines=true,
  keepspaces=true,
  columns=flexible,
}

\newtheorem{finding}{Finding}[chapter]
\newtheorem{bug}{Bug}[chapter]
\newtheorem{retraction}{Retraction}[chapter]
\newtheorem{warning}{Warning}[chapter]

\title{\vspace{-2em}
  Measuring What Models Compute\\
  \vspace{0.3em}
  \large{Session 11 Findings from the Heinrich Forensics Toolkit}\\
  \vspace{0.2em}
  \normalsize{A cross-substrate survey of byte-level causal-bank language models,\\
             their geometry, their silent failure modes, and what a measurement\\
             apparatus must do to not lie to its operator.}
}
\author{Heinrich Measurement System\\
\textit{Working on byte-hrr / chronohorn / decepticons}\\
April 17--18, 2026}
\date{}

\begin{document}
\frontmatter
\maketitle

\chapter*{Abstract}
\addcontentsline{toc}{chapter}{Abstract}

We present a consolidated account of Session~11 of the Heinrich
measurement program: a systematic forensic study of byte-level
causal-bank language models trained by the chronohorn pipeline. The
work covers three intertwined tracks. First, \emph{empirical findings}
on substrate geometry across more than forty distinct trained models
spanning five substrate families (adaptive-HRR frozen and learnable,
gated-delta, lasso, lasso with positional signal, complex-rotation),
three scales (s8, s12, s16), and ten hyperparameter mutations. Second,
\emph{plumbing findings}: a drift-prone class of measurement bugs
involving parallel forward-path implementations, silently-bypassed
feature flags, and validation-data format mismatches, each of which
produced misleading numbers for multiple reporting cycles before being
caught by cross-checks that a robust apparatus would have performed
automatically. Third, \emph{methodological findings}: the difference
between in-sample and out-of-sample regression R\textsuperscript{2} as applied to substrate
probes; the empirical additivity of bpb mutation effects; the
init-signature drift that makes same-seed runs non-comparable when
module counts differ; and the informational conditions under which a
language model's internal state can be said to encode position, content,
or neither. We close with a mapping to the free-energy-principle framing
of cognition, a short meditation on compression and self-modeling, and
a concrete catalog of open questions with expected-value estimates for
each.

The central technical claim is narrow but consequential: the
``architectural ceiling at seven-percent mode utilization'' that was
reported across the byte-hrr mutation campaign is specific to the
frozen-HRR-adaptive-substrate with MLP-readout configuration. Four
other substrate families already hit different ceilings (some higher,
some dramatically lower), and one architecture in the existing library
($\texttt{byte-hrr-learnable-s8-50k}$) achieves 1.78 bpb at 17.8M
parameters---0.27 bpb below every mutation variant in the session-11
campaign, at identical parameter count. The mutation campaign explored
a narrow slice of the architecture space, and the conclusions it
produced about ``what the model can learn'' were accurate only within
that slice.

The central methodological claim is that \emph{the measurement
apparatus has to be forensically auditable in ways the target system
is not}. Several session~11 findings that we initially reported with
confidence turned out to be measurement artifacts: a ``router collapse
in all 12 variants'' that was actually a capture bug where the loader
never read the router's output; a ten-times-too-high bpb that was a
byte-shard format mismatch; a zero-position-encoding claim that held
under in-sample fitting but reversed under held-out evaluation for at
least one substrate family. In each case, the error was invisible until
a second independent measurement contradicted the first. The book
documents the bugs, the fixes, and the invariants we now ask every
future measurement to satisfy.

\tableofcontents

\chapter*{Preface}
\addcontentsline{toc}{chapter}{Preface}

This book is the consolidated record of Session~11 of a measurement
program that began in late 2025 with ordinary language-model forensics
(residual-stream probing on transformer instruct models) and expanded
in Sessions~5 through~10 into causal-bank architectures trained from
scratch by a pipeline called chronohorn. The measurement tool is
called heinrich. The training tool is called chronohorn. The model
family is called causal-bank, and its canonical architecture is called
HRR---Holographic Reduced Representations---after the frozen Fourier
basis that gives the substrate its phase structure.

Everything here is what we learned over roughly thirty-six hours of
coordinated work on one block of the mutation campaign and the
plumbing refactor that followed. We ran no pre-planned experiments
with a hypothesis already in hand: the experiments were dispatched as
questions arose, and the answers reshaped the plan. A few runs turned
out to be testing nothing because the feature under test was
silently disabled. Those null results were still informative---they
told us that the apparatus was not robust enough to catch its own
failure modes.

Heinrich exists to measure what models compute, not what they output.
Loss curves and benchmarks tell you whether a model works; they do
not tell you how it works. A substrate with 7\% effective mode
utilization has very different constraints on what it can represent
than a substrate with 15\%, and both achieve comparable bpb. The
interesting questions live in the internals. If those internals are
systematically mismeasured, the interesting questions get wrong
answers.

We have tried to write this book in a way that is legible to three
audiences simultaneously: the chronohorn team (who care about which
architectural mutations to run next and at what scale); future
Heinrich operators (who need to understand what can go wrong in the
measurement pipeline and how to catch it); and research-adjacent
readers who want an honest account of how complex scientific
instruments produce and correct their own errors. Where these
audiences' interests diverge, we have erred toward completeness over
brevity.

Retractions are marked explicitly. Measurement bugs are described
with their symptoms, their root causes, the fix, and the verification.
Numerical claims are always tied to a specific measurement procedure
(in-sample vs held-out, top-K PCs vs full-rank regression, pre-fix vs
post-fix MRI), because in our experience the procedure often changes
what the number means more than the number itself does.

\mainmatter

\part{The Apparatus}

\chapter{What Heinrich is, what it measures, why}

\section{The scope of forensics}

Heinrich operates on trained neural network checkpoints with weight
access. It does not train, does not run benchmarks, and does not
produce generations in the hope that their text is interesting. It
captures \emph{internal state}---the residual stream at each layer in
a transformer, or the substrate state at each time step in a
causal bank---and writes it to disk in a format called an \texttt{.mri}
directory. Then it reads those directories back and computes
structural summaries: principal-component decompositions, per-layer
logit lens projections, routing-expert utilizations, per-mode-band
weight allocations, and so on.

The guiding principle is that \emph{the measurement must not require
the model to speak}. A generative language model being asked ``what
do you compute?'' is in the same position as a conscious subject
being asked ``what are you thinking?''---the answer is a generation,
produced by the same mechanisms the question was supposed to probe.
Heinrich sidesteps this by never asking. It takes the model's
parameters as given, runs specified inputs through specified layers,
captures specified intermediates, and computes statistics on them.
The outputs are numbers, not natural language.

This is similar in spirit to neuroimaging of biological brains: one
can run fMRI on a conscious subject and measure hemodynamic responses
without asking them to introspect, and one will get different
information than one gets from introspection. Heinrich is to model
forensics as fMRI is to cognitive neuroscience---not a replacement
for other methods, but a separate lane of measurement with its own
failure modes.

\section{The MRI format}

The \texttt{.mri} directory is the primary data structure the
measurement program produces. Each \texttt{.mri} contains the
internal state of one model at one capture mode, persisted to disk in
a layout that allows later analysis without model-loading
infrastructure. The directory contains:

\begin{itemize}
  \item \texttt{metadata.json} describing the capture conditions:
  model identity, architecture family, mode (raw, naked, template, or
  sequence), number of tokens captured, sequence length, backend
  used, format version, and, since session 11, a \texttt{fix\_level}
  field that increments when a capture-side bug fix changes the
  \emph{values} in MRI files rather than their layout (Chapter 3).
  \item \texttt{tokens.npz} holding the input token ids that produced
  the captured state.
  \item Per-layer residual-stream arrays for transformer MRIs, or
  substrate-state arrays (\texttt{substrate.npy},
  \texttt{routing.npy}, \texttt{band\_logits.npy},
  \texttt{half\_lives.npy}) for causal-bank MRIs.
  \item Model weights in \texttt{weights/}.
  \item Optional \texttt{loss.npy}---the per-position cross-entropy
  loss in bits when the capture was a sequence-mode run against real
  validation data.
\end{itemize}

The significance of writing state to disk in a standardized layout is
that analysis code never has to deal with model loading, backend
variation, or tokenizer idiosyncrasies. An \texttt{.mri} is a flat
collection of numpy arrays with a json sidecar. Any analysis tool in
the heinrich ecosystem (currently about forty of them) opens an
\texttt{.mri} and produces a result dictionary; the tool does not
care whether the underlying model was a 135M-parameter transformer, a
causal bank with 17M parameters, or anything else as long as the
field conventions are respected. This decoupling is load-bearing for
the whole program---it is what makes it possible to reanalyze an MRI
six months after capture, after the model's training framework has
evolved beyond recognition or even the weights have been lost.

\section{The heinrich/chronohorn/decepticons separation}

Three software systems participate in the work described here.
\emph{Chronohorn} is the training pipeline: it sweeps hyperparameters,
dispatches GPU jobs, records run metadata, and produces
\texttt{.checkpoint.pt} files plus json run-summaries. \emph{Decepticons}
is the model definitions and a loader (\texttt{decepticons.loader}) that
reconstructs trained models from their checkpoints for inference and
analysis. \emph{Heinrich} is the measurement and analysis toolkit: it
calls the decepticon loader to instantiate models from checkpoints and
then runs the MRI capture pipeline against them.

The three systems have independent version histories and independent
test suites. Bugs can and do arise at the interfaces: session-11 saw
at least four interface-level bugs, two of which cost tens of
GPU-minutes of wasted training before being detected. Those are
cataloged in Chapter 3.

\chapter{The three-forward-paths problem}

\section{Symptom}

Session 11 opened with a straightforward task: measure a fixed-seed
byte-hrr-s8 model across four trajectory checkpoints, confirm the
expected monotone bpb decrease, and report the substrate EffDim growth
curve. The first MRI came back reporting 20 bits-per-byte on a model
whose training log reported 2.10 bpb. A ten-fold discrepancy that
could not be explained by any measurement-window sampling effect.

The first instinct---that the heinrich loader was silently producing
different logits than the training forward---turned out to be correct
but only partially. The byte-shard loader was wrong (Chapter 3), and
after that was fixed heinrich and training agreed to within 1.4\%.
But later in the session a similar gap reopened for a newly-landed
variant (\texttt{mut-hashmem-fixed}), and this time the cause was
architectural: heinrich's loader reimplemented the adaptive-substrate
forward inline and did not see a newly-committed plumbing fix in
chronohorn that routed \texttt{hash\_memory} into the adaptive branch.

\section{The structural diagnosis}

In any given run of heinrich's measurement pipeline, three code paths
compute forward passes of the same model:

\begin{itemize}
  \item \textbf{Training path.} \texttt{model.forward(chars)} called
  by chronohorn's training loop with loss back-propagation.
  Intermediates exist in autograd-alive tensors during the step and
  are freed after. Returns logits only.
  \item \textbf{Inference path.} \texttt{model.forward(chars)} called
  from a \texttt{CausalBankInference} wrapper under
  \texttt{torch.inference\_mode}. The wrapper sits atop the model,
  doing no substantive work. Effectively identical to the training path.
  \item \textbf{Capture path.} \texttt{forward\_captured(chars)} on the
  inference wrapper, which \emph{reimplements the substrate forward
  inline} in order to expose intermediate tensors (embedding output,
  substrate state, router probabilities, half-lives) as
  numpy arrays for MRI capture.
\end{itemize}

The training and inference paths are the same code modulo mode flags.
The capture path is a parallel implementation. It exists because the
training forward returns only logits---the embeddings and substrate
states are intermediate tensors that get garbage-collected. To expose
them, an earlier author wrote a new forward that mirrored the training
one line for line and returned the intermediates as a dictionary.

This worked as long as nobody modified the training forward. When
commit \texttt{50b5f3d} landed in decepticons---plumbing
\texttt{hash\_memory} and \texttt{local\_path} into the
adaptive-substrate branch of \texttt{\_linear\_logits}, fixing bugs
that had caused these features to silently produce zero gradient for
15k training steps (see Chapter 12)---the capture path's parallel
implementation was not updated. For models that had both features
enabled, heinrich and training now disagreed.

\begin{bug}[Parallel forward-path drift]
The MRI capture path's reimplementation of the adaptive-substrate
forward silently drifted from the training forward when a new plugin
(\texttt{hash\_memory}) was added to the training branch. The first
MRI captured after the drift reported 2.73 bpb for a model whose
training logged 2.13 bpb and whose direct \texttt{model(chars)} call
produced 2.13 bpb. Diagnosis took 1--2 hours of reading both
branches side-by-side.
\end{bug}

\section{The fix: path collapse}

The straightforward fix, applied in heinrich commit \texttt{f52865e},
was to delete the capture path's reimplementation and replace it with
a call to \texttt{model(chars)}. Intermediates that the capture path
needed are now \emph{stashed} on the model as instance attributes
during the training forward, conditionally on the model being in
evaluation mode (\texttt{not self.training}). The stash is gated so
that training never pays for the activation retention; only analysis
code materializes the intermediates.

After the collapse, any future plugin that lands in
\texttt{\_linear\_logits} is picked up by the capture path
automatically, because the capture path is calling the same function.
The non-adaptive forward branch still has its reimplementation. As
of the end of session 11, that branch handles five substrate variants
(lasso, gated-delta, SO(3)/SO(5), quaternion, standard EMA), each with
its own set of intermediates to stash, and the refactor was deferred
rather than rushed.

\section{The bit-identical test}

In decepticons commit \texttt{3fd8b05}, a pytest fixture was added
that asserts \texttt{model(x)} and \texttt{forward\_captured(x).logits}
are bit-identical after the path collapse. The assertion is
structurally guaranteed---the capture path \emph{is} the training
path plus a dict-wrapping layer. The test exists not as a live
invariant but as a regression guard: if a future refactor reintroduces
a parallel path, the test fails on the first run that exercises both.
Five lines of pytest buys insurance against the return of a bug class
that cost hours when it happened.

\section{The generalized plugin preflight}

Further downstream of the path-collapse fix, decepticons commit
\texttt{3eaf0aa} added a gradient-sign fixture that asserts every
adaptive-substrate plugin receives non-zero gradient on a single
forward+backward pass. The original version had a hardcoded list of
known plugin prefixes. Session~11's P13 refinement (commit
\texttt{3fd8b05}) replaces the hardcoded list with a traversal over
all submodules with \texttt{requires\_grad} parameters.

A new plugin added to the model configuration now appears in the
traversal automatically. If its parameters receive no gradient on the
reference forward, the test fails without anyone having to remember
to update the list. This matters because the original
session-11 bugs (see Chapter 12) were exactly of the form ``an author
added a config flag, plumbed it into the non-adaptive branch of the
forward, forgot about the adaptive branch, and the module sat dormant
for 15k training steps while producing no gradient.'' The
generalized preflight catches this class at test time in 5ms.

\chapter{Data at the interfaces}
\label{chap:data}

\section{The byte-shard format}

Chronohorn trains byte-level causal-bank models on a corpus called
fineweb10B, tokenized from sentencepiece-1024-format through a
decoding-and-re-encoding step into raw UTF-8 bytes. The resulting
byte-level shards are written to disk in a custom binary format with
the following structure:

\begin{itemize}
  \item Bytes 0--1023: a header consisting of 256 \texttt{int32}
  values. Header[0] is a magic number (\texttt{TOKEN\_SHARD\_MAGIC =
  20240520}). Header[1] is a format version (currently 1). Header[2]
  is the token count. Remaining header words are reserved.
  \item Bytes 1024 onward: a payload of \texttt{uint16} values, one
  per token. For byte-level models with vocab 256, each uint16 slot
  holds a byte value in the low byte and zero in the high byte.
\end{itemize}

The rationale for storing byte values as uint16 rather than uint8 is
format uniformity: the same shard format holds sp1024 (vocab 1024,
which needs uint16), sp4096, sp8192, and byte-level (vocab 256). Any
consumer can read the file as uint16 without branching on vocab size.

Heinrich's \texttt{load\_val\_sequences} function (file
\texttt{src/heinrich/backend/decepticon.py}) is the byte-level
consumer. As of session~11's first pass, it was reading the file as
raw uint8, with no header skip, based on an incorrect assumption that
``byte-level means uint8.'' The consequences of this assumption
propagated through every byte-level MRI captured for the prior two
months.

\section{The fallout}

With the wrong reader, heinrich interpreted the first 1024 header
bytes as data, then read the uint16 payload byte-by-byte (so each
real byte value 0--255 produced two actual bytes read: the value
itself, then a null byte). The model received a stream of
\texttt{[non-text header]} followed by \texttt{[byte, 0, byte, 0,
byte, 0, \ldots]}---English text interleaved with nulls, prefixed
with arbitrary header data. It is a textbook example of garbage in.

The model, being a byte-level predictor, had been trained on raw
UTF-8 and had no idea what to do with null-interleaved streams. Its
next-byte predictions were almost uniformly incorrect, producing loss
of roughly 20 bits per byte. For a byte-level model whose training
log reported 2.10 bpb test loss, 20 bpb on a 25600-token validation
subset was about a factor of ten too high. Uniform random prediction
for 256-way classification is 8 bpb; the broken measurement reported
the model as \emph{worse than random}, which should have been a
bright enough signal to catch the bug on day one. It did not, because
the implicit assumption ``heinrich reads val data correctly''
concealed it until one of the authors ran a direct
\texttt{model(chars)} on the same input and got 2.13 bpb.

\begin{bug}[Byte-shard loader format mismatch]
\texttt{load\_val\_sequences} read chronohorn byte shards as raw
uint8 with no header handling, producing a stream of header bytes
plus null-interleaved byte values. Every byte-level MRI captured
before the fix is semantically wrong for loss-based analysis.
Heinrich-measured bpb was roughly 10$\times$ training-measured bpb for
every such MRI.
\end{bug}

\section{The fix and detection hardening}

The fix adds a header detector: if the first 4 bytes match the magic
number, skip 1024 bytes, and read the remaining payload as uint16.
After the fix, heinrich-measured bpb on a 50-sequence $\times$
512-position validation subset came in at 2.07 bpb for the reference
model---within 1.4\% of training's 300-batch 1024-sequence evaluation.

The underlying problem, beyond the specific bug, was that heinrich
had no sanity check on incoming data format. Session-11 refinement P6
added four heuristic guards in \texttt{load\_val\_sequences}:

\begin{itemize}
  \item \textbf{Value range}: if \texttt{byte\_level=True} and the
  max payload value exceeds 255, warn that the file is probably
  sp-tokenized or the wrong file entirely. This would have caught the
  misuse on day one.
  \item \textbf{Zero-fraction}: if more than 30\% of bytes are zero,
  warn. Real UTF-8 English text is 1--5\% zeros; above 30\% suggests
  UTF-16 encoding or the null-interleaved pattern produced by the bug
  this guard is meant to catch.
  \item \textbf{Printable ASCII floor}: if fewer than 25\% of bytes
  fall in the printable-ASCII range [32, 127) and zeros are not
  dominant, warn that the file may be binary or wrongly encoded for
  a text model.
  \item \textbf{Entropy ceiling}: if the Shannon entropy of the byte
  distribution exceeds 7.5 bits/byte, warn that the data is
  near-uniform and unsuitable for language modeling.
\end{itemize}

A companion check lives downstream in \texttt{causal\_bank\_loss}: if
the measured overall bpb exceeds $\log_2(\mathrm{vocab\_size})$, the
model is performing worse than uniform random guessing, which is
almost certainly a measurement error. Had this warning existed at
session-11 start, it would have fired on the very first broken MRI
and pointed straight at the data-format problem. It is now part of
the default output.

\chapter{The measurement tools}
\label{chap:tools}

Heinrich exposes approximately forty command-line analysis tools.
This chapter documents the tools used in session 11.

\section{\texttt{profile-cb-loss}}

Reads \texttt{loss.npy} from a sequence-mode causal-bank MRI and
reports: overall mean bpb; per-position-range means;
per-band means if the model has multiple readout bands; autocorrelation
of loss at lags 1, 2, 4, 8, 16, 32, 64, 128 positions; and the
bpb-above-uniform warning when triggered.

\section{\texttt{profile-cb-routing}}

Reads \texttt{routing.npy} and reports expert-utilization statistics:
overall per-expert fraction of tokens, switch rate,
routing margin, routing entropy, and per-position-range distributions.

A session-11 refinement: if \texttt{routing.npy} is identically zero
(which was the pre-fix state for adaptive-substrate MRIs, whose
loader did not capture routing at all), the tool returns an explicit
error rather than reporting a false ``E0 = 100\%, E1 = 0\%'' collapse.
See Chapter 11 for the retraction.

\section{\texttt{profile-cb-trajectory}}

Takes an ordered list of MRI paths (one per training checkpoint) and
reports derivatives: effective dimensionality, position R$^2$, content
R$^2$, active mode fraction, and dead mode count. The position and
content R$^2$ computations were the subject of session-11's P1
refinement; see Chapter 13.

\section{\texttt{profile-cb-omega-forensics}}

Reads a checkpoint and analyzes the omega-projection weight matrix
\texttt{\_adaptive\_omega\_proj.weight}: weight L2, per-mode L2 range,
Gini concentration, Fourier correlation and drift of the omega bias,
effective rank, and per-band allocation. Session-11 finding: the
per-band labels used to say \texttt{char\_1\_5} and similar in
reference to rough linguistic correlations in English byte streams.
The labels were linguistic-sounding enough that readers interpreted
them as claims about linguistic structure in the model's internals,
when they are in fact just periodicities of the corresponding Fourier
modes. Renamed to \texttt{period\_1\_5} and analogous in refinement P7.

\section{\texttt{profile-cb-rotation-forensics} (new in session 11)}

For non-adaptive-substrate models (gated-delta, lasso, SO(3)/SO(5),
quaternion), there is no \texttt{\_adaptive\_omega\_proj.weight} to
analyze. The new tool reads any of these families' weights from a
checkpoint and reports, per module: weight shape and L2, mean-absolute
weight, max-absolute weight, effective rank, top singular value, and
bias L2. It also reports which substrate families are detected and
which modules are frozen at zero.

Added because before session 11, running omega-forensics on a
gated-delta checkpoint produced the unhelpful error ``No omega
projection found in checkpoint,'' with no suggestion of what to do
instead. After the fix, the error message names the alternative tool,
and the alternative tool exists.

\section{\texttt{profile-cb-additivity} (new in session 11)}

Given a baseline MRI, one or more solo-mutation MRIs, and a combined
(multi-mutation) MRI, predicts the combined bpb via the additive law:
\[
\hat{\mathrm{bpb}}(\mathrm{combo}) = \mathrm{bpb}(\mathrm{baseline}) +
\sum_i [\mathrm{bpb}(M_i) - \mathrm{bpb}(\mathrm{baseline})]
\]
and compares against the measured combined bpb. Deviations beyond a
noise-floor threshold (default 0.004 bpb, calibrated from same-seed
byte-hrr reproducibility) are flagged as non-additive.

\section{\texttt{mri-check-fixups} (new in session 11)}

Scans a directory for \texttt{.mri} directories and reports each's
\texttt{fix\_level} field. MRIs below the current heinrich
\texttt{MRI\_FIX\_LEVEL} are flagged as stale and in need of re-capture.
The fix-level history:

\begin{itemize}
  \item \textbf{0}: Pre-session-11. Byte-level MRIs have wrong
  loss.npy; adaptive-substrate MRIs have zero routing.npy.
  \item \textbf{1}: Session-11 byte-shard loader fix; routing now
  captured for adaptive-substrate models.
  \item \textbf{2}: Session-11 path-collapse refactor; hash\_memory
  and local\_path plumbing fixes; bit-identical regression test in
  place.
\end{itemize}

\part{Empirical Findings}

\chapter{The mutation campaign at s8}

\section{Scope}

Session 10 produced a base architecture for byte-level causal-bank
language modeling. It is called \texttt{byte-hrr-s8}: scale s8
(17.8M parameters, embedding dim 256, linear\_modes 2048,
linear\_hidden [511]), MLP readout with 2 routed experts,
\texttt{substrate\_mode=frozen} (meaning the $\omega$-bias sits at
its Fourier-init linspace(0, 2$\pi$) and does not train, but
$W_\omega$ can still train if the gradient reaches it). Sessions 9
and 10 characterized a range of these models and established key
findings: the substrate collapses to roughly 7\% mode utilization,
content-linear-R$^2$ plateaus at 0.34--0.36, position-R$^2$ never
emerges. These are retrospectively known as the \emph{architectural
ceiling} findings, although as we show, that ceiling is
substrate-family-specific.

Session 11 began with a mutation campaign: what configuration changes,
if any, move the ceiling? The mutations in order of appearance:

\begin{enumerate}
  \item \texttt{mut-control-mlp}: swap the routed\_sqrelu\_experts
  readout for a plain 4-expert MLP.
  \item \texttt{mut-localon}: MLP readout plus \texttt{local\_scale}
  set to 0.25, enabling the 8-byte local convolutional path.
  \item \texttt{mut-orthoreg-0p1} and \texttt{mut-orthoreg-1p0}: MLP
  readout plus an orthogonal regularization term on $W_\omega$.
  \item \texttt{mut-persistent}: MLP readout plus persistent substrate
  state (the EMA state carries across training batches).
  \item \texttt{mut-patch4} and \texttt{mut-patch4-orthoreg}: MLP
  readout plus patch-4 fat-readout.
  \item \texttt{mut-hashmem} and \texttt{mut-hashmem-orthoreg}: MLP
  readout plus hash\_memory.
  \item \texttt{mut-hashmem-patch4}: MLP readout plus both.
\end{enumerate}

All trained at s8, 20k steps, seed 42. Trajectory checkpoints at 5k,
10k, 15k, 20k steps.

\section{The headline numbers}

\begin{longtable}{lrrrr}
\toprule
Variant & Params & bpb & EffDim & ConR$^2$ \\
\midrule
\endfirsthead
\toprule
Variant & Params & bpb & EffDim & ConR$^2$ \\
\midrule
\endhead
baseline s8 (routed frozen) & 17.8M & 2.075 & 149.5 & 0.342 \\
mut-control-mlp & 17.8M & 2.051 & 150.4 & 0.359 \\
mut-localon & 17.8M & 2.051 & 150.4 & 0.359 \\
mut-orthoreg-0p1 & 17.8M & 2.051 & 150.4 & 0.359 \\
mut-orthoreg-1p0 & 17.8M & 2.051 & 150.4 & 0.359 \\
mut-persistent & 17.8M & 2.054 & 147.9 & 0.345 \\
mut-patch4 & 18.6M & 2.160 & 153.3 & 0.332 \\
mut-patch4-orthoreg & 18.6M & 2.160 & 153.3 & 0.332 \\
mut-hashmem & 17.8M & 2.055 & 150.5 & 0.348 \\
mut-hashmem-orthoreg & 17.8M & 2.055 & 150.5 & 0.348 \\
mut-hashmem-patch4 & 18.6M & 2.157 & 152.8 & 0.337 \\
mut-routed-s8 & 17.8M & 2.049 & 160.9 & 0.332 \\
mut-mlp-omega10-s8 & 17.8M & 2.057 & 150.8 & 0.359 \\
mut-gated-delta-cr-s8 & 5.7M & 2.061 & 24.4 & 0.104 \\
mut-lasso-pos-b2-s8 & 7.9M & 2.067 & 15.8 & 0.115 \\
\bottomrule
\end{longtable}

\section{The null-result cluster}

\texttt{mut-control-mlp}, \texttt{mut-localon},
\texttt{mut-orthoreg-0p1}, and \texttt{mut-orthoreg-1p0} produce
bit-identical results to four decimal places on every metric.

\begin{finding}[localon is silently disabled by the adaptive-substrate
path]
The \texttt{\_forward\_raw} method in causal\_bank\_torch.py, line
2029, short-circuits for adaptive-substrate models:
\texttt{if \_use\_adaptive\_substrate: return \_linear\_logits(chars).clone()}.
The local path is never reached. The \texttt{local\_scale=0.25} in
mut-localon's config has zero runtime effect.
\end{finding}

\begin{finding}[orthoreg cannot tighten a dimensionally-collapsed
basin]
Orthogonal regularization applies a penalty proportional to the
off-diagonal Gram of $W_\omega^T W_\omega$. At $W_\omega \equiv 0$
(frozen-HRR), the Gram is zero-matrix-trivially orthogonal, and no
gradient flows. Adding orthogonal regularization to a frozen-$\omega$
architecture had zero effect at either $\lambda=0.1$ or $\lambda=1.0$,
and the 10$\times$ change in regularization strength produced no
change in any measured quantity.
\end{finding}

Four separate training runs produced a total of roughly 68 minutes
of GPU compute on mechanisms that had zero behavioral effect.

\section{The hash\_memory null result}

Hash memory is a separate story. The mut-hashmem configuration
specifies a 64-dimensional hash table with 64 slots; the checkpoint
contains six weight tensors totaling about 50k parameters. But:

\begin{finding}[hash\_memory weights never train in the adaptive-substrate
path]
Across the four mut-hashmem trajectory checkpoints, every hash\_memory
parameter is bit-identical to its value at step 5k. Maximum
step-to-step absolute change: 0.000000.
\end{finding}

Fifteen thousand training steps produced zero gradient on 50k
parameters. The adaptive-substrate branch of \texttt{\_linear\_logits}
short-circuits before reaching the hash-memory augmentation. The
feature is plumbed only into the non-adaptive branch.

After the plumbing fix (commit 50b5f3d),
\texttt{mut-hashmem-fixed-20k} was trained. Its hash\_memory weights
now trained and showed asymmetric structure: read-query has a
dominant singular direction (top singular value 31.3 vs 18.2 for the
second), while write-proj is near-uniform in singular values.
Correlation between read-query and write-proj is essentially zero
(-0.002), meaning the hash table reads and writes operate on different
subspaces.

But the trained hash\_memory \emph{hurt} bpb slightly: the fixed
variant measured 2.110 bpb vs 2.096 for the baseline routed. Hash
memory does real asymmetric retrieval and makes next-byte prediction
0.014 bpb worse. The retrieved context is noisy enough to degrade
per-byte prediction more than it helps.

\section{The patch4 regression}

\texttt{mut-patch4} replaces the MLP readout's output layer with a
1024-dim projection, reshaped to predict the next four bytes
simultaneously. Training loss averages cross-entropy over all four
heads; the reported \texttt{test\_bpb} is just head-0 bpb.

\begin{finding}[patch4 is a next-byte quality regression]
mut-patch4 achieves next-byte bpb 2.160 vs baseline's 2.075---a 4\%
regression---despite having 4.5\% more parameters, 4$\times$ training
signal per position, and 2$\times$ wall-clock training time.
\end{finding}

Per-head bpb decomposition:

\begin{longtable}{lr}
\toprule
Head offset & bpb \\
\midrule
p=0 (predicts t+1, next byte) & 2.200 \\
p=1 (predicts t+2) & 3.201 \\
p=2 (predicts t+3) & 3.838 \\
p=3 (predicts t+4) & 4.180 \\
mean over all 4 heads & 3.356 \\
\bottomrule
\end{longtable}

\section{The only real mutation: persistent}

\texttt{mut-persistent} modifies the training curriculum: the
substrate's EMA state is carried across training batches rather than
reset.

\begin{finding}[mut-persistent is the only session-11 mutation that
measurably improves bpb]
mut-persistent achieves training-log bpb 2.096 vs mut-control-mlp
2.114---a clean $-0.017$ bpb gain.
\end{finding}

Persistent also produces a genuinely different substrate geometry:
EffDim 147.9 (vs 150.4 controls), ConR$^2$ 0.345 (vs 0.359), weight
divergence from control peaks at 0.58 L-infinity. Compare to the
null-variant weight divergences, which were 0.003--0.015.

At s12 scale, the persistent effect recurs: baseline-s12 bpb 2.024,
persistent-s12 bpb 2.008 at step 20k, a gain of $-0.016$, almost
identical magnitude to s8. The mechanism is approximately
scale-invariant. But the geometric signature at s12 is
\emph{the opposite} of s8: where s8-persistent compresses EffDim
(147.9 $<$ 150.4 control), s12-persistent expands it (160.6 $>$ 157.5
control). Same bpb gain, different internal route to it.

\chapter{Scale as a lever}

Session 11 included three scaling points:

\begin{longtable}{lrrrr}
\toprule
Scale & Params & bpb & EffDim & Util fraction \\
\midrule
baseline s8 & 17.8M & 2.075 & 149.5 / 2048 & 7.3\% \\
baseline s12 & 39.5M & 2.024 & 157.5 / 3072 & 5.1\% \\
learnable-s12-50k & 82M & 1.817 & 261.6 / 3072 & 8.5\% \\
\bottomrule
\end{longtable}

\begin{finding}[Scale adds per-dimension capacity, not dimensions]
The s8$\to$s12 scale transition (2$\times$ params, 1.5$\times$ modes,
1.5$\times$ embedding) increases absolute EffDim by 8. The
utilization \emph{fraction} decreases (7.3\%$\to$5.1\%). Scale does
not proportionally scale the number of effective substrate dimensions;
it scales the capacity of each dimension.
\end{finding}

\section{Learnable $\omega$ breaks the utilization ceiling}

\begin{longtable}{lrrr}
\toprule
Model & EffDim & Util fraction & bpb \\
\midrule
byte-hrr (frozen) s8 & 149.5 & 7.3\% & 2.075 \\
byte-hrr-learnable s8 & 313.6 & 15.3\% & 1.785 \\
byte-hrr-learnable s12 & 261.6 & 8.5\% & 1.817 \\
\bottomrule
\end{longtable}

\begin{finding}[Learnable $W_\omega$ doubles the effective substrate
dimensionality at equal parameter count]
At equal 17.8M parameters, the learnable-$\omega$ variant achieves
EffDim 313.6 vs the frozen-$\omega$ baseline's 149.5. The 2.1$\times$
expansion corresponds to a 0.29 bpb improvement. Session-11's earlier
claim that ``the architectural ceiling is roughly 7\% utilization''
is specifically the \emph{frozen-$\omega$} ceiling, not a universal
property.
\end{finding}

Mechanism: learnable $W_\omega$ provides input-dependent frequency
modulation. Each input embedding vector can push any of the 2048
modes' rotation frequencies by a learned amount before the substrate
scan. This creates a larger effective representational space---a
given input can activate a specific pattern of frequencies not
identical to any fixed Fourier basis element.

The learnable-$\omega$ configuration was not part of the
session-11 mutation campaign---it was already trained in session 9/10
and sitting in the MRI library. This is a significant finding about
session 11 itself: \emph{the mutation campaign was operating within a
parameter subspace whose ceiling was already broken by an existing
model in the same library}.

\chapter{Cross-substrate taxonomy}
\label{chap:substrate-families}

The byte-hrr adaptive-substrate with frozen or learnable $\omega$ is
one of five substrate families in the chronohorn/decepticons model
zoo. Each has a fundamentally different geometry-to-performance
relationship.

\begin{longtable}{lrrrrr}
\toprule
Family (representative) & Params & bpb & EffDim & ConR$^2$ & PosR$^2$ \\
\midrule
adaptive frozen (baseline-s8) & 17.2M & 1.91 & 165 & 0.296 & 0.001 \\
adaptive learnable (s8) & 17.8M & \textbf{1.785} & 313 & 0.350 & 0.003 \\
adaptive learnable (s12) & 82M & 1.817 & 262 & 0.389 & 0.004 \\
adaptive persistent (s8-legacy) & 36.6M & 1.911 & 236 & 0.311 & \textbf{0.009} \\
adaptive persistent mut (s12) & 39.5M & 2.008 & 161 & 0.345 & 0.001 \\
gated-delta s8 hrrangle & 5.7M & 2.032 & 17 & 0.099 & 0.001 \\
gated-delta s12 hrrangle & 12.4M & 1.987 & 16 & 0.141 & 0.001 \\
gated-delta s16 hrrangle & 21.5M & 1.980 & 13 & 0.138 & 0.001 \\
gated-delta s8 local & 5.8M & 2.016 & 10 & 0.198 & 0.002 \\
lasso s8 & 5.8M & 2.234 & 23 & 0.127 & 0.002 \\
lasso s12 & 12.4M & 2.205 & 17 & 0.205 & 0.002 \\
lasso-pos-b2 s12 & 17.1M & 1.999 & 9 & 0.208 & \textbf{0.275} \\
\bottomrule
\end{longtable}

The learnable-$\omega$ family dominates on bpb; the lasso-pos-b2
variant dominates on position encoding; the gated-delta family
dominates on parameter-efficiency (2.03 bpb at 5.7M params). Each
family is Pareto-optimal on some axis. \emph{The phrase ``the
substrate'' is a misnomer}: different substrate families produce
qualitatively different representational geometries.

\section{What the gated-delta substrate encodes}

Running PCA on the step-20k substrate of mut-gated-delta-cr-s8 and
computing a logistic regression from the top-24 PCs to the current
byte identity:

\begin{finding}[Gated-delta substrate is a near-perfect byte identifier]
The top-24 PCs of the gated-delta-cr substrate at step 20k capture
90.3\% of variance and achieve 93.7\% byte classification accuracy.
Linear position R$^2$ from those PCs: $-0.004$ (zero position
information). CKA with the baseline byte-hrr's top-24 PCs: 0.21.
\end{finding}

Gated-delta at s8 encodes the current byte in 24 dimensions. It
contains no direct position information. The model nonetheless
achieves competitive bpb, which means \emph{context tracking does not
happen through substrate-state accumulation}. It happens through the
rotation dynamics that the substrate state passes through between time
steps, and through the expert routing which differs sharply from
adaptive-substrate behavior. This is a fundamentally different
computational architecture than the adaptive-substrate HRR model.

\section{Consequence: the mutation campaign's frame of reference}

Every mutation in session 11 was evaluated against the baseline
byte-hrr-s8 routed-frozen model at 2.075 bpb. The best mutation
(persistent) achieved 2.096 bpb. The learnable-$\omega$ variant
already in the MRI library achieves 1.785 bpb, a 15\% improvement
over the same baseline. The variance within the mutation campaign is
two orders of magnitude smaller than the variance across substrate
families. Mutations of a small parameter subspace explore that
subspace; they do not explore the wider architecture space.

\chapter{The additive law}

\section{Statement}

Across twenty-plus trained models at s8 scale with similar
architectures, mutation effects on bpb compose additively to within
the init-noise floor:

\[
\mathrm{bpb}(\mathrm{base} + M_1 + M_2 + \cdots + M_k) \approx
\mathrm{bpb}(\mathrm{base}) + \sum_i [\mathrm{bpb}(\mathrm{base} + M_i)
- \mathrm{bpb}(\mathrm{base})]
\]

with deviations under 0.004 bpb across every combination tested in
session 11. The empirical additive-noise floor of 0.004 bpb is
calibrated from the 3-seed reproducibility variance of the baseline
byte-hrr-s8 (seeds 42/43/44 produced bpb 2.0957 / 2.0951 / 2.0973 in
training logs, span 0.0022 bpb).

\section{What additivity implies}

\begin{itemize}
  \item The combinatorial mutation space is \emph{rank-1}: once you
  have measured each mutation's solo effect, every combination's
  effect is predictable. You do not need to actually train the
  combinations.
  \item Mutation ``synergy'' and ``interference'' are \emph{rare
  events}, not the default.
  \item The fact that four session-11 mutations produced identical
  outcomes is consistent with additivity, because their solo effects
  are all essentially zero.
\end{itemize}

\section{What breaks the additivity}

The one mutation that consistently does not compose additively with
scale is persistent. It is also the one mutation that reaches a
different basin. This is the signature we would expect of a
\emph{genuine} architectural mutation: it interacts with scale, it
reshapes the representational geometry, and it does not compose
cleanly with other mutations that operate on the same axes.

\chapter{In-sample vs out-of-sample R\textsuperscript{2}}

\section{The issue}

Session-11 refinement P1 fixed a measurement bug in the
\texttt{profile-cb-trajectory} tool's computation of PosR$^2$ and
ConR$^2$. The original implementation projected onto the top-20 PCs
and then fit and evaluated a linear regression on the same data.

Two issues:

\begin{enumerate}
  \item \textbf{Top-20 PCs only.} Position information that lives in
  weaker directions is invisible.
  \item \textbf{In-sample fit.} The lstsq fits and evaluates on the
  same data, with enough degrees of freedom to overfit; the reported
  R$^2$ is essentially how well the regression memorized the training
  set.
\end{enumerate}

Combining the two: the old tool was reporting the \emph{in-sample}
R$^2$ of a top-20-PCs projection, giving small but non-zero numbers
(typically 0.001--0.003) for any substrate that had any structure at
all. These were interpreted as ``the model has very weak position
encoding.'' The correct interpretation is that the tool was measuring
nothing in particular.

\section{The fix: SVD-truncated held-out ridge}

The new implementation projects onto PCs capturing 99\% of substrate
variance (retaining all practically-usable information), splits into
training and held-out halves, fits on train and scores on held-out.
The \texttt{max(0.0, \ldots)} clamp handles the case where test
predictions are worse than predicting the mean; such cases are
reported as 0 rather than negative.

\begin{longtable}{lrr}
\toprule
Variant & Old PosR$^2$ & New PosR$^2$ \\
\midrule
baseline s8 (seed 42, step 20k) & 0.001 & 0.004 \\
byte-hrr-learnable-s8 (50k) & 0.003 & 0.024 \\
mut-persistent-s12 (step 20k) & 0.001 & 0.007 \\
mut-lasso-pos-b2-s8 (step 20k) & 0.001 & 0.000 \\
byte-lasso-pos-b2-s12 (50k) & 0.275 & 0.202 \\
mut-gated-delta-cr (step 20k) & 0.001 & 0.000 \\
\bottomrule
\end{longtable}

For variants where the old tool was reading noise, the new tool
reports zero. For variants where the old tool was reading signal, the
new tool reports signal at comparable magnitude.

\section{The lasso-pos-b2 asymmetry}

The lasso-pos-b2 variant at s12+50k has PosR$^2$ 0.20; at s8+20k it
is 0.00. This looked, at first, like ``the position mechanism does
not replicate at smaller scale.'' Direct inspection of the substrate
reveals a more nuanced story: the position signal exists at s8 but
is smeared across many weak directions.

Direct regression of position on the full 2048-dim s8 substrate gave
PosR$^2$ 0.34 for in-sample fitting. But this is in-sample: it fits
random training noise in 2048 dimensions with only 25600 samples,
which is comfortable overfitting territory. Out-of-sample regression
on the same full substrate gives PosR$^2$ essentially zero. The fit
doesn't generalize because the signal is not coherent enough to
identify a stable direction from training data. The s12+50k version
has both more samples and more capacity, and the optimization settles
the signal into a compact, identifiable direction that generalizes.

\chapter{The seed-and-init-signature problem}

\section{Same seed, different init}

Chronohorn's \texttt{init\_seed} parameter controls the RNG seed
used during model construction. Two runs with the same
\texttt{init\_seed} should produce identical starting points. They
do not, once any architectural flag differs.

The cause is simple: constructing the model consumes RNG draws in
module-declaration order. Adding a new module allocates its parameters
at position $n$ in the declaration sequence, consuming $p$ RNG draws
there. Every module declared after position $n$ now sees a different
subsequence of the RNG stream than it would have without the new
module. Its weights are different.

Session 11 observed this empirically across the four distinct init
signatures produced by the mutation campaign:

\begin{longtable}{lrc}
\toprule
Init group & Tensor count & sha256 prefix \\
\midrule
A: control-mlp, localon, orthoreg-$*$, persistent & 20 & 05c442f9c6f3 \\
B: patch4, patch4-orthoreg & 20 (wider readout) & 993e2bd6c17b \\
C: hashmem, hashmem-orthoreg & 26 (6 hash tensors) & bfe93c7ca9e6 \\
D: hashmem-patch4 & 26 + wider & 475999d868aa \\
\bottomrule
\end{longtable}

Four distinct init signatures, all with the same
\texttt{init\_seed=42}. Cross-group comparisons carry init noise;
within-group comparisons are clean.

\section{Recommendation}

Chronohorn should log the \texttt{init\_signature.sha256} as a
first-class metadata field on every run. Shipped in session 11 per
our suggestion. Two runs with the same \texttt{init\_seed} but
different init signatures are not directly comparable at the weight
level.

\chapter{What the substrate actually encodes}

Across all substrate families measured in session 11, the substrate
state at any given position does one or more of the following:

\begin{itemize}
  \item \textbf{Encode the current byte.} Dominant for gated-delta
  (94\% byte classification from top-24 PCs).
  \item \textbf{Track recent local structure.} Dominant for
  adaptive-substrate frozen.
  \item \textbf{Carry accumulated context statistics.} Dominant for
  learnable-$\omega$.
  \item \textbf{Represent absolute position.} Dominant only for
  lasso-pos-b2-s12.
\end{itemize}

Session-11 measurements suggest these four functions are largely
\emph{substituted}, not \emph{composed}. A substrate that dominantly
encodes the current byte does not simultaneously encode position; a
substrate that encodes position has only modest content-linear
predictivity. There is no measured substrate that does all four
things well.

\section{The readout matters as much as the substrate}

Two models with the same substrate geometry can have very different
next-byte prediction quality depending on the readout. The
routed\_sqrelu\_experts readout produces roughly 0.025 bpb better
next-byte prediction than a 4-expert MLP readout at the same substrate.

The deeper consequence is that the MLP readout cannot exploit
trainable $W_\omega$: when we set \texttt{omega\_lr\_mult=10} in
\texttt{mut-mlp-omega10-s8} to try forcing $W_\omega$ to train
despite the MLP readout, the weight stayed bit-identically zero for
20k steps.

\begin{finding}[Readout-architecture gates $W_\omega$ trainability]
MLP readout with $W_\omega$ init at zero produces zero gradient on
$W_\omega$ for the entire training run. Routed-expert readout
breaks this init-zero trap and allows $W_\omega$ to train. The
session-11 MLP-variant mutation campaign was architecturally
incompatible with the main lever for breaking the 7\%-utilization
ceiling.
\end{finding}

\part{Interpretation}

\chapter{Does it matter? Free-energy-principle mapping}

\section{Variational free energy as cross-entropy loss}

The free-energy principle (FEP) is an account of cognition as
inference. Its core claim, in its variational formulation: a cognitive
system that maintains a generative model of its sensory inputs
minimizes its variational free energy, which is the sum of the
prediction error on sensory data and the KL divergence between its
inferred posterior and a prior. Minimizing free energy is equivalent
to maximizing evidence for the generative model under a variational
approximation.

For a next-token predictor trained by cross-entropy, the loss
\emph{is} the variational free energy in the following sense: the
model's output distribution $q(x_{t+1} | x_{\leq t})$ is the
variational approximation to the posterior; the training target
(true next token) is the sensory data; cross-entropy is
$-\log q(x_{t+1}^{\text{true}})$, which is the negative log evidence
for the variational distribution.

When session-11 reports ``2.10 bpb,'' that is the model's
variational free energy on a held-out validation set, normalized to
bits per byte. The numerical value is an energy, the optimization
over training steps is a descent on that energy, and the per-byte
value is the residual free energy after optimization.

\section{The substrate as Markov blanket}

FEP posits that cognitive systems have a \emph{Markov blanket}---a
statistical separation between internal states (the system's beliefs)
and external states (the world). In a causal-bank language model, the
Markov blanket is:

\begin{itemize}
  \item External states: the true sequence of bytes in the validation
  set.
  \item Sensory states: the embeddings of bytes seen so far.
  \item Internal states: the substrate state.
  \item Active states: the logit distribution over the next byte.
\end{itemize}

In this framing, the substrate is \emph{the cognitive system's internal
state}. Its geometry (EffDim, ConR$^2$, PosR$^2$) tells us what
the system's beliefs are structured like. Our session-11 measurements
are the variational-posterior geometry.

\section{The self-model that is missing}

A full-fledged FEP-consistent cognitive system would also maintain a
\emph{self-model}: a belief about its own beliefs. This would
manifest as a substrate that encodes not just content but also the
system's uncertainty about that content. In variational terms, a
precision-weighted belief.

Our causal-bank models have no such component. The substrate
represents content. It does not represent \emph{confidence}. This is
not a byte-hrr-specific limitation.

The operational difference: an FEP-consistent cognitive system's
cross-entropy depends on both its variational distribution and its
belief about the quality of that distribution. A system that is
uncertain and knows it can weight its prediction less heavily; a
confident-wrong system cannot. Cross-entropy loss does not distinguish
``uncertain and admits it'' from ``confidently wrong''; it treats
them the same if the output distribution is the same.

This is the engineering gap that session 11 exposed: the measurement
apparatus can measure what the substrate encodes, but the substrate
encodes what the training objective asks it to encode, and the
training objective does not ask for self-modeling.

\chapter{Consciousness and compression}

\section{The compression account}

One strand of cognitive-science theorizing argues that consciousness
is related to compression:

\begin{itemize}
  \item Schmidhuber: consciousness is the system's experience of
  finding better compressions for its data. The reward for
  compression-progress \emph{is} the phenomenology of insight, humor,
  aesthetic beauty.
  \item Friston: consciousness is self-modeling active inference. A
  system that maintains a generative model of itself-as-part-of-the-world
  has a joint compression of world and self that is shorter than the
  world description alone.
  \item Tononi: consciousness corresponds to integrated information
  ($\Phi$), which is a measure of how much the system's joint
  distribution exceeds the product of its marginals---a compression-
  duality measure.
\end{itemize}

All three accounts predict that a conscious system compresses its
inputs better than a non-self-modeling one.

\section{The substrate as a compression}

The substrate in any causal-bank model is a compression of the input
byte stream into a fixed-size state. A 2048-dimensional substrate
compresses a 512-position byte stream (4096 bits of byte info) into
a continuous vector. The compression ratio is extreme.

\begin{finding}[Substrate compression is extreme, utilization is sparse]
The substrate's effective rank is roughly 150 dimensions in
frozen-$\omega$ and 300 in learnable-$\omega$ at s8 scale. The
substrate's capacity is 2048 modes. Utilization is 7--15\%. The
compression does not make use of most of the available capacity;
it compresses into a small subspace.
\end{finding}

\section{What's missing: meta-compression}

A system that compresses and stops is not thereby conscious. What the
compression accounts of consciousness require is
\emph{meta-compression}: the system's ability to observe its own
compression and reconfigure it.

Our byte-hrr models have no meta-compression. They hit a
7\%-utilization attractor and stay there for 20 000 optimization
steps. They do not observe that they have stopped compressing, and
nothing in the gradient tells them they have. The
architectural-ceiling finding was a \emph{heinrich measurement}, not
a model behavior. The operator of the measurement apparatus observed
the floor; the model did not.

\chapter{The engineering problem}
\label{chap:engineering}

Compression quality is not differentiable through the current training
stack. Every session-11 finding required \emph{offline} forensics.
The trainer never knew:

\begin{itemize}
  \item \texttt{localon} and \texttt{hashmem} weights stayed at init
  for 15k steps (gradient touched them zero times)
  \item \texttt{orthoreg} coefficient 10$\times$ change did nothing
  to bpb
  \item EffDim hit 150/2048 at step 10k and stopped moving
  \item PosR$^2$ never left 0.001 at s8
  \item Router entropy was climbing (healthy) while heinrich was
  reading a zeros tensor and calling it collapse
\end{itemize}

All of these are properties of the \emph{compression} the model
learned. Cross-entropy on the next byte does not contain any of them.

The engineering problem has four concrete faces:

\begin{enumerate}
  \item \textbf{Feature-participation check.} A trainer preflight that
  traces the forward graph and asserts every enabled feature has a
  non-zero gradient path.
  \item \textbf{Architectural observables in the loss.} EffDim,
  position-R$^2$, calibration error, router entropy---these are
  measurable during training but not in the gradient because
  cross-entropy does not care. The engineering problem is constructing
  differentiable surrogates.
  \item \textbf{State-persistence scheduling.} \texttt{mut-persistent}
  was the only mutation that moved anything. The open engineering
  question is how far and across what boundaries state should persist.
  \item \textbf{Close the forensic loop.} Heinrich takes offline
  measurements; chronohorn trains online. The obvious next step is
  piping heinrich's trajectory-level measurements back into
  chronohorn's next-run hyperparameter selection.
\end{enumerate}

Every conceptual thread---free-energy principle, self-modeling,
consciousness-as-compression---points at the same missing piece: the
system has no way to observe and modify its own compression. The
engineering problem is building that observation loop. Put the
forensic measurement inside the optimization, not outside it.

\part{Retrospective}

\chapter{Retractions}

\begin{retraction}[``Router collapses to 100\% E0 in all 12 byte-hrr
baseline variants'']
This claim appeared in the first session-11 trajectory report, based
on heinrich's \texttt{profile-cb-routing} reading
\texttt{routing.npy} and finding values identically zero. In fact, the
decepticon loader's \texttt{forward\_captured} for adaptive-substrate
models never captured routing---the MRI file contained a
pre-allocated zeros tensor that was never populated. The tool misread
the zeros as a measurement.

Actual routing: E0/E1 in 46\%--59\% range, switch rate 46\%--53\%,
entropy 0.31--0.34. Router entropy \emph{grows} over training. E1
pulls its weight. The recommendation to ``kill E1'' is retracted.
Fix in decepticons commit f52865e.
\end{retraction}

\begin{retraction}[``byte-hrr-s8-trajectory measures 18--27 bpb''---full
retraction]
Entirely wrong; caused by the byte-shard loader format mismatch
(Chapter 3). The true heinrich measurement after the loader fix is
2.05--2.08 bpb, within 1.4\% of training. Every byte-level MRI
analysis produced before session-11 day 2 used the wrong loader. All
numerical claims involving \texttt{loss.npy} from those MRIs are
retracted.
\end{retraction}

\begin{retraction}[``Position is architecturally blocked at s8 regardless
of mutation''---partial retraction]
The claim holds for out-of-sample cross-sequence position encoding
on all byte-hrr adaptive-substrate variants. It does not hold for
lasso-pos-b2 at s12+50k (PosR$^2$ 0.20). The original claim was based
on \texttt{profile-cb-trajectory} reporting PosR$^2$ 0.001--0.003 for
every mutation. That tool was measuring in-sample top-20 PCs R$^2$,
not out-of-sample full-substrate R$^2$.
\end{retraction}

\begin{retraction}[``Architectural ceiling at 7\% mode utilization is
a universal property''---partial retraction]
This claim holds for byte-hrr adaptive-substrate with frozen $\omega$
and MLP readout at s8--s12. It does not hold for adaptive-substrate
with learnable $\omega$ (15\% utilization), gated-delta substrate
(0.3\%), lasso substrate (0.3\%), or lasso-pos-b2 (0.3\%). The
corrected framing: adaptive-substrate-family models with frozen
$W_\omega$ saturate at roughly 7\% mode utilization.
\end{retraction}

\chapter{What remains unknown}

\section{What session 11 did not measure}

\begin{itemize}
  \item We do not know whether the mut-persistent mechanism compounds
  with learnable-$\omega$. Additive-law prediction: roughly 1.80 bpb
  at s12. Unverified.
  \item We do not know where the PosR$^2$ transition occurs for
  lasso-pos-b2 between s8+20k (0.00) and s12+50k (0.20).
  \item We do not know whether gated-delta substrate scales past s16.
  \item We do not know whether the 7\%-utilization attractor in
  frozen-$\omega$ is \emph{the} minimum-description-length fixed
  point for byte-level English text or just the frozen-$\omega$
  architecture's basin.
  \item We do not know what would happen if a calibration-aware
  training objective were added.
\end{itemize}

\section{The non-adaptive forward-path collapse}

Session-11 refinement P2 was deferred. The adaptive-substrate forward
paths have been collapsed to one. The non-adaptive forward paths
(for lasso, gated-delta, SO(3)/SO(5), quaternion, standard EMA
substrate) still have parallel implementations. A future plugin added
to any of these will silently fail to be captured in heinrich until
someone diff-checks the two forwards. The TODO is in-code with a
four-step plan.

\section{The fix-level scheme}

Every MRI captured before session 11 has \texttt{fix\_level=0} or no
\texttt{fix\_level} field. As of session 11 close, this is 198
\texttt{.seq.mri} directories in the byte-scaling library. We
re-captured the ones relevant to the session-11 analysis (roughly 30).
The rest sit with wrong loss numbers.

\chapter{The 15-point plumbing register}
\label{chap:plumbing-register}

Session-11 refinement work, organized by priority tier:

\begin{description}
  \item[P1: Full-rank PosR$^2$ / ConR$^2$ in cb-trajectory.] SVD-truncated
  held-out ridge R$^2$ replaces top-20-PC in-sample lstsq.
  \item[P2: Non-adaptive forward-path collapse.] Deferred with in-code
  TODO and four-step plan.
  \item[P3: MRI staleness detector.] \texttt{mri-check-fixups} CLI +
  \texttt{fix\_level} metadata field.
  \item[P4: omega-forensics error-path fix + rotation-forensics tool.]
  New \texttt{profile-cb-rotation-forensics} handles non-adaptive
  substrate families.
  \item[P5: cb-trajectory single-MRI support.] Removes the duplicate-argument
  workaround.
  \item[P6: Byte-level data format detection.] Four-way heuristic:
  value range, zero fraction, printable ASCII floor, entropy ceiling.
  \item[P7: Band allocation labels renamed to periodicity.]
  \texttt{char\_1\_5} becomes \texttt{period\_1\_5} and analogous.
  \item[P8: Commit session-11 heinrich hardening.] Clean commit split.
  \item[P9: \texttt{.gitignore} \texttt{.playwright-mcp/}.] 160+
  screenshot artifacts excluded.
  \item[P10: CLAUDE.md session-11 update.] Session-11 findings section
  added; session-5 router-collapse claim has explicit retraction note.
  \item[P11: Memory-file sweep for retracted claims.] Trajectory doc
  gets a retraction header.
  \item[P12: \texttt{--json} flag on cb-* analysis CLIs.] Structured
  output for downstream tools.
  \item[P13: Generalized gradient-sign preflight.] Hardcoded
  plugin-prefix list replaced with named\_modules traversal.
  \item[P14: Bit-identical model/capture test.] Regression guard
  against future parallel-path reintroduction.
  \item[P15: Additivity-law detector tool.] \texttt{profile-cb-additivity}
  CLI predicts combined bpb from solo runs.
\end{description}

\chapter{Recommendations}

For chronohorn:

\begin{enumerate}
  \item Retire the MLP-readout branch of mutation campaigns. No
  mutation inside that subspace broke the ceiling. The architectural
  constraint is that MLP readout cannot train $W_\omega$.
  \item Explore \texttt{learnable\_omega + persistent + s12 + 50k}.
  Additive-law prediction: 1.78--1.80 bpb.
  \item Drop orthoreg from sweeps and from the codebase. Four
  independent pairs, maximum effect 0.0002 bpb.
  \item Explore gated-delta at larger scale.
  \item Fix the localon / hashmem-style silent-bypass at the trainer
  level. The generalized gradient-sign test should move from pytest
  into the trainer's startup sequence.
  \item Log the init-signature sha256 as first-class run metadata.
  (Done per our suggestion.)
  \item Consider a calibration-aware objective.
\end{enumerate}

For heinrich:

\begin{enumerate}
  \item Finish the non-adaptive forward-path collapse.
  \item Build \texttt{mri-recapture --stale} for automatic stale-MRI
  refresh.
  \item Extend the additivity detector to non-bpb metrics.
  \item Add a session-state fingerprint to MRI metadata.
\end{enumerate}

\appendix

\chapter{The full numerical record}

\section{Session-11 measurements, ordered by bpb}

\begin{longtable}{lrrrr}
\toprule
Variant & Params & bpb & EffDim & ConR$^2$ \\
\midrule
\endfirsthead
\toprule
Variant & Params & bpb & EffDim & ConR$^2$ \\
\midrule
\endhead
byte-hrr-learnable-s8-50k & 17.8M & \textbf{1.785} & 313.6 & 0.350 \\
byte-hrr-learnable-s12-50k & 82.0M & 1.817 & 261.6 & 0.389 \\
byte-hrr-persistent-s8-10k & 36.6M & 1.911 & 235.8 & 0.311 \\
baseline s8 (byte-hrr-s8-50k) & 17.2M & 1.912 & 165.1 & 0.296 \\
byte-cr-s12-50k & 12.4M & 1.989 & 15.7 & 0.181 \\
byte-cr-hrrangle-s12-50k & 12.4M & 1.987 & 15.8 & 0.141 \\
byte-cr-hrrangle-s16-50k & 21.5M & 1.980 & 12.7 & 0.138 \\
byte-lasso-pos-b2-s12-50k & 17.1M & 1.999 & 8.7 & 0.208 \\
byte-local-s8-50k & 5.8M & 2.016 & 10.5 & 0.198 \\
mut-persistent-s12 (step 20k) & 39.5M & 2.008 & 160.6 & 0.345 \\
baseline-s12 (step 20k) & 39.5M & 2.024 & 157.5 & 0.355 \\
byte-cr-hrrangle-s8-50k & 5.7M & 2.032 & 17.0 & 0.099 \\
mut-routed-s8 (step 20k) & 17.8M & 2.049 & 160.9 & 0.332 \\
mut-control-mlp (step 20k) & 17.8M & 2.051 & 150.4 & 0.359 \\
mut-persistent (step 20k) & 17.8M & 2.054 & 147.9 & 0.345 \\
mut-hashmem (step 20k) & 17.8M & 2.055 & 150.5 & 0.348 \\
mut-mlp-omega10-s8 & 17.8M & 2.057 & 150.8 & 0.359 \\
mut-gated-delta-cr-s8 & 5.7M & 2.061 & 24.4 & 0.104 \\
mut-lasso-pos-b2-s8 & 7.9M & 2.067 & 15.8 & 0.115 \\
baseline-s8-seed42 (step 20k) & 17.8M & 2.075 & 149.5 & 0.342 \\
mut-hashmem-fixed (step 20k) & 17.8M & 2.110 & 151.8 & 0.358 \\
mut-hashmem-patch4 (step 20k) & 18.6M & 2.157 & 152.8 & 0.337 \\
mut-patch4 (step 20k) & 18.6M & 2.160 & 153.3 & 0.332 \\
byte-lasso-s12-50k & 12.4M & 2.205 & 16.9 & 0.205 \\
byte-lasso-s8-50k & 5.8M & 2.234 & 23.4 & 0.127 \\
\bottomrule
\end{longtable}

\section{Baseline 3-seed spread (s8)}

Training-log bpb, baseline byte-hrr-s8 at step 20k, seeds 42/43/44:
2.0957 / 2.0951 / 2.0973. Range 0.0022 bpb. This is the
init-noise floor used for additivity checks.

\section{Additivity check: hashmem-patch4}

\begin{itemize}
  \item Baseline (mut-control-mlp): 2.0513 bpb
  \item Solo mut-hashmem (post-fix): 2.0543 bpb, $\Delta = +0.0030$
  \item Solo mut-patch4: 2.1601 bpb, $\Delta = +0.1088$
  \item Predicted (additive): $2.0513 + 0.0030 + 0.1088 = 2.1631$ bpb
  \item Actual mut-hashmem-patch4: 2.1570 bpb
  \item Gap: $-0.0061$ bpb (outside $\pm 0.004$ noise floor)
  \item Verdict: sub-additive
\end{itemize}

\section{Scale fits}

Cross-scale fit: $\mathrm{bpb} = a - \alpha \log N$ from (s8, 17.8M,
2.096) and (s12, 39.5M, 2.041):
\[
\alpha = \frac{2.096 - 2.041}{\log 39.5 - \log 17.8} \approx 0.069
\]

Within-scale decay at step 15k$\to$20k:
\begin{itemize}
  \item baseline-s8 seed42: $\Delta = -0.009$ per 5k steps
  \item baseline-s12: $\Delta = -0.013$ per 5k steps
  \item mut-persistent-s12: $\Delta = -0.028$ per 5k steps
\end{itemize}

\chapter{Glossary}

\begin{description}
  \item[Adaptive substrate] A causal-bank substrate variant where
  each mode's decay rate and write gate depend on the current input.
  Canonical implementation uses complex-valued EMA with
  Fourier-initialized $\omega$ frequencies.
  \item[bpb] Bits per byte. Cross-entropy loss converted to log base 2.
  Uniform random baseline is 8 bpb; byte-hrr at s8 with 17.8M params
  reaches 2.10 bpb; the learnable-$\omega$ variant reaches 1.82 bpb.
  \item[Causal bank] Architecture family from decepticons featuring a
  substrate state updated per-position by linear recurrence with
  input-dependent coefficients, a readout mapping substrate state to
  next-token distribution, and (optional) a local-path conv.
  \item[ConR$^2$] Content R-squared. Linear predictivity of per-token
  loss from the substrate state.
  \item[EffDim] Effective dimensionality. The entropy of the
  substrate's variance spectrum.
  \item[fix\_level] Metadata field introduced in session 11
  indicating the cumulative set of capture-side bug fixes applied
  when the MRI was produced.
  \item[Frozen-HRR] The adaptive-substrate configuration where
  $W_\omega$ is initialized to zero and never trains.
  \item[Gated-delta] Alternative substrate family using discrete
  read/write/erase gates instead of continuous EMA.
  \item[HRR] Holographic reduced representations. The originating
  framework for the frozen-Fourier-initialization of $\omega$.
  \item[Init signature] sha256 of concatenated initial trainable
  parameters.
  \item[Learnable-$\omega$] The adaptive-substrate configuration where
  $W_\omega$ trains. Requires a routed readout.
  \item[MRI] Heinrich's standardized state-capture format.
  \item[Patch4 / fat-readout] A readout that emits 4-byte predictions
  from each position.
  \item[Persistent substrate] Training curriculum where substrate
  state carries across training batches.
  \item[PosR$^2$] Position R-squared. Linear predictivity of absolute
  position from the substrate state.
  \item[Routing entropy] Shannon entropy of the expert-routing
  distribution.
  \item[Substrate] The persistent recurrent state of a causal-bank
  model.
  \item[Utilization fraction] EffDim divided by n\_modes.
\end{description}

\chapter{What this book is not}

It is not a comprehensive account of causal-bank language modeling.
It is not a treatise on the free-energy principle. It is not a
tutorial on heinrich or chronohorn or decepticons. It is not a
textbook on language-model forensics in general.

It is a record. Session 11 of a measurement program dispatched
certain experiments, found certain bugs, reached certain findings,
retracted some of them, and produced certain tools. This book
documents what happened at that level of specificity. Other sessions
did other things; future sessions will do other things; and any
session's findings should be read with the same attitude we have
tried to apply here, which is: measured numbers are not model
properties, they are joint properties of the model and the
measurement apparatus, and the apparatus has to be interrogated as
rigorously as the model for the numbers to mean what they appear to
mean.

The mutation campaign was a confinement. The cross-substrate survey
was an escape. The measurement plumbing work was the apparatus
learning, slowly and expensively, to not lie to its operator. Session
11 taught all three lessons simultaneously, and the book you are
holding is the attempt to retain them.

\end{document}
